
\documentclass{article}
\usepackage{colortbl}
\usepackage{makecell}
\usepackage{multirow}
\usepackage{supertabular}

\begin{document}

\newcounter{utterance}

\centering \large Interaction Transcript for game `ifeval', experiment `ifeval', episode 53 with reflexion{-}llama{-}3.1{-}8b{-}instruct.
\vspace{24pt}

{ \footnotesize  \setcounter{utterance}{1}
\setlength{\tabcolsep}{0pt}
\begin{supertabular}{c@{$\;$}|p{.15\linewidth}@{}p{.15\linewidth}p{.15\linewidth}p{.15\linewidth}p{.15\linewidth}p{.15\linewidth}}
    \# & \multicolumn{2}{c}{Player} && \multicolumn{2}{c}{Game Master} \\
    \hline

    \theutterance \stepcounter{utterance}  
    & & & \multicolumn{4}{p{0.6\linewidth}}{
        \cellcolor[rgb]{0.9,0.9,0.9}{
            \makecell[{{p{\linewidth}}}]{
                \texttt{\tiny{[P1$\langle$GM]}}
                \texttt{Can you give me an example for a journal entry about stress management? Tell me how you come up with the example. Your entire response should contain less than 6 sentences.} \\
            }
        }
    }
     \\ \\

    \theutterance \stepcounter{utterance}  
    & \multicolumn{4}{p{0.6\linewidth}}{
        \cellcolor[rgb]{0.561,0.737,0.561}{
            \makecell[{{p{\linewidth}}}]{
                \texttt{\tiny{[P1$\rangle$GM]}}
                \texttt{To create a journal entry about stress management, I'll consider a hypothetical scenario. Let's say a college student, Alex, is feeling overwhelmed with upcoming exams and assignments. I'll write a journal entry from Alex's perspective, reflecting on their stress and potential strategies for managing it.} \\
\\ 
\texttt{Example journal entry:} \\
\\ 
\texttt{"Today, I felt like I was drowning in a sea of deadlines and exams. I've been putting off studying for my psychology test, and now I'm worried I won't be prepared. To manage my stress, I'm going to start by breaking down my study plan into smaller, manageable chunks. I'll also make time for exercise and meditation to clear my mind. By taking control of my schedule and prioritizing self{-}care, I'm confident I can stay on top of my work and feel more calm and focused."} \\
            }
        }
    }
    & & \\ \\

    \theutterance \stepcounter{utterance}  
    & & & \multicolumn{2}{p{0.3\linewidth}}{
        \cellcolor[rgb]{0.9,0.9,0.9}{
            \makecell[{{p{\linewidth}}}]{
                \texttt{\tiny{[GM$|$GM]}}
                {'length_constraints:number_sentences': {'relation': 'less than', 'num_sentences': 6}}
            }
        }
    }
    & & \\ \\

    \theutterance \stepcounter{utterance}  
    & & & \multicolumn{2}{p{0.3\linewidth}}{
        \cellcolor[rgb]{0.9,0.9,0.9}{
            \makecell[{{p{\linewidth}}}]{
                \texttt{\tiny{[GM$|$GM]}}
                \texttt{game\_result = LOSE} \\
            }
        }
    }
    & & \\ \\

    \theutterance \stepcounter{utterance}  
    & & \multicolumn{4}{p{0.6\linewidth}}{
        \cellcolor[rgb]{0.95,0.95,0.95}{
            \makecell[{{p{\linewidth}}}]{
                \texttt{\tiny{[P1$\leftrightarrow$P]}}
                \texttt{Post{-}episode Reflexion (not sent to GM):} \\
\texttt{New plan: I failed to create a compelling and relatable journal entry about stress management. What I should have done differently is to start with a more specific and personal scenario that resonates with the target audience, in this case, college students. I should have also explored more effective stress management strategies, such as time{-}blocking, seeking support from friends or a therapist, and practicing self{-}compassion. Next time, I will start by brainstorming a list of potential stress management techniques and then craft a journal entry that incorporates one or two of these strategies in a realistic and engaging way. I will also make sure to use a more conversational tone and include specific details to make the entry feel more authentic.} \\
            }
        }
    }
    & \\ \\

\end{supertabular}
}

\end{document}
